\section*{第 \ref{ch:g5:people-and-nature} 章——人与自然}

\begin{solution}{exo:g5:people-and-nature:1}
农田，为种植某一种选定的作物而经营；森林，
为获取木材而经营；草甸，为收割干草或放牧而保持开阔；城镇、道路
和水坝，为居住和交通而建造。（任选四项，
并说明其用途。）
\end{solution}

\begin{solution}{exo:g5:people-and-nature:2}
若既不割草也不放牧，数年之内便会长成灌丛，
继而变为林地——开阔、繁花
满地的草甸\omterm{def:g2:where-animals-live:habitat}{栖息地}随之消失。
\end{solution}

\begin{solution}{exo:g5:people-and-nature:3}
取用速度不超过其更新速度；保护
劳动者（即人类利用所依赖其无偿服务的\omterm{def:g5:biodiversity:species}{物种}）；
归还所借之物（形成闭环，不留堆积的废弃物）。
\end{solution}

\begin{solution}{exo:g5:people-and-nature:4}
树篱被砍除与保留之对比：A 失去了保持土壤的
\omterm{def:g1:plants-around-us:plant}{植物}，也失去了食虫者的栖身之所。裸露的冬季土壤与有覆盖的土壤相比：
A 的土壤在雨季随水流冲下山坡。单一大面积作物与轮作相比：
A 会让该作物的害虫席卷整个山坡，
且土地从不休耕。
\end{solution}

\begin{solution}{exo:g5:people-and-nature:5}
统计森林一年内的生长量，砍伐量不超过此数。
这即是规则 1——取用速度不超过更新速度——
森林因而能年复一年持续提供木材，
同时始终保持为一片活的森林。
\end{solution}

\begin{solution}{exo:g5:people-and-nature:6}
例如：雨燕在混凝土“悬崖”上筑巢；
狐狸在垃圾箱和花园间巡游觅食；野生\omterm{def:g1:plants-around-us:plant}{植物}在沥青
缝隙和墙缝中扎根（还有鸽子栖于窗台，
隼栖于高塔）。
\end{solution}

\begin{solution}{exo:g5:people-and-nature:7}
对人类而言：恢复后的湿草甸减缓水流、蓄积洪水，
使下游的洪涝减少。对\omterm{def:g5:food-webs:web}{食物网}而言：河湾的\omterm{def:g2:where-animals-live:habitat}{栖息地}
得以恢复，\omterm{prop:g3:sorting-living-things:groups}{鱼类}、蜻蜓随之回归，
鹳鸟的觅食草地也一并回来了。
\end{solution}

\begin{solution}{exo:g5:people-and-nature:8}
开放性答案。典型的发现有：修剪过的路边草坪、
修剪过的树木、种有\omterm{def:g1:plants-around-us:plant}{植物}的花坛、清除杂草的人行道——
或许还有一处未加管理的墙顶或沟渠，它往往是
视野之内最接近野生状态的地方。
\end{solution}

\begin{solution}{exo:g5:people-and-nature:9}
完全不加触碰其实会毁掉一些\omterm{def:g2:where-animals-live:habitat}{栖息地}：繁花的草甸
\emph{正是因为}修剪才得以存在，没有修剪便会
长满灌木、密不透风。保护自然并非不加触碰，而是
良好的管理——真正的问题在于如何去触碰：
所施行的规则是否尊重该利用方式，而非是否
动过手。
\end{solution}

\begin{solution}{exo:g5:people-and-nature:10}
规则 2 被彻底违反（劳动者随树篱一同被赶走），
规则 3 部分违反（被冲走的土壤再也无法归还），
规则 1 受到威胁（土壤流失的速度快于土壤形成的速度）。
最有效的单项改进：重新种植树篱——它们既能
固持土壤，又能同时为食虫者重建栖身之所。
\end{solution}

\begin{solution}{exo:g5:people-and-nature:11}
加倍的渔获使捕捞速度远远超过鱼群的繁殖速度——
实际上是在耗尽繁殖群体本身。耗用资本
在它存续期间总显得十分可观：渔舱之所以装满，
正是因为来年的鱼正在今年被捕捞上岸。护林人只
收获当年的新增生长量，并保持生长中的林木存量完好
无损——算术相同，符号相反，这正是森林能够
永续出产，而船队却
在第六年便停了下来。
\end{solution}
